%==============================================================================
%  sections_v3/proofs_core_lemmas.tex -- proofs of Lemmas 1-5 of
%  sections_v3/theorem_section.tex (card IDs D1, L1, L2, L3, L4), in the
%  ledger's own order.
%
%  SOURCE OF TRUTH: research/model_v4/MODEL_CARD.md, version stamp 2026-08-30
%  (F5 route R40-A, gate PASS, theory-record freeze;
%  commit 65b8db3), section 6 result ledger.  Each proof
%  below proves the statement exactly as theorem_section.tex carries it --
%  that is, the amended card row.  Where a proof file on record predates an
%  amendment, the amendment is what is proved and the repair the card names is
%  the one incorporated.
%
%  This file states nothing: every result is stated in theorem_section.tex and
%  every hypothesis in model_section.tex.  No environment is defined here; the
%  only environment used is amsthm's proof.  Proof-local notation follows the
%  card's 4.6 rulings.
%==============================================================================

\begin{proof}[Proof of Lemma~\ref{lem:D1} (D1)]
Throughout this proof, $j(\cdot)$ denotes the cutoff selection map of
Definition~\ref{def:cpbe}\,(i), the map carrying the signal into a plan at the cutoff vector in
force, and $\sigma(\cdot)$ the $\sigma$-algebra generated by its argument. Parts~(a) and~(b) are established first, under no restriction whatever on
the probability of either cell; part~(c) follows.

\emph{Step 1 (one probability space, and a Borel selection map).} By Assumption~\ref{asm:A1} the
vector $(v,\varepsilon,\xi,z_{0:H})$ has a joint law on a finite product of Polish spaces, the noise
coordinates taking values in the finite set $\{-\bar z,0,+\bar z\}$ of
Definition~\ref{def:primitives}, so every object below lives on one probability space. The signal
$s = v+\varepsilon$ is a sum of two coordinates and is Borel. By Definition~\ref{def:cpbe}\,(i),
$j(\cdot)$ is a step function with breakpoints $k_1 \le \dots \le k_{J-1}$, finitely many because the
plan menu is finite by Assumption~\ref{asm:A2p}. A step function with finitely many breakpoints is a
finite sum of indicators of half-lines, so $j(\cdot)$ is Borel and $\{s : j(s) = j\}$ is Borel for
every $j \in \mathcal J$.

\emph{Step 2 (the engagement flag is measurable).} The realised flag $a = a_{j(s)}$ composes the
Borel map of Step~1 with the map $j \mapsto a_j$ of Definition~\ref{def:plans}, defined on the finite
set $\mathcal J$. A map on a finite set carrying the discrete $\sigma$-algebra is measurable, and
$\mathcal J$ is finite by Assumption~\ref{asm:A2p}; hence $\{a=1\}$ is Borel and lies in $\sigma(s)$.

\emph{Step 3 (the crossing date is attained, unique and measurable on the Voice region).} Fix $j$
with $a_j = 1$. By clause (TR-ii) of Assumption~\ref{asm:TR}, $s \mapsto B_j(s,d)$ is weakly
increasing for Voice plans, and a weakly monotone real function is Borel because the preimage of
$[\tau,\infty)$ is an interval; so $\{s : B_j(s,d) \ge \tau\}$ is Borel at each date $d$. The
calendar $\{0,\dots,H\}$ is finite by Assumption~\ref{asm:A2p}, so
$\{c_j \le d\} = \bigcup_{d' \le d}\{s : B_j(s,d') \ge \tau\}$ is a finite union of Borel sets, and
$c_j$ is a Borel map into $\{0,\dots,H\} \cup \{+\infty\}$. Its value is unique: a nonempty subset of
the finite well-ordered set $\{0,\dots,H\}$ has exactly one minimum, and Assumption~\ref{asm:A4}
makes $c$ that first date; if the subset is empty then $c_j = +\infty$ by Definition~\ref{def:plans}.

\emph{Step 4 (the disclosure indicator is measurable).} By Definition~\ref{def:plans}, $D=1$ requires
$a=1$, so
\begin{equation}
\label{eq:D1-cellunion}
\{D = 1\} \;=\; \bigcup_{j \in \mathcal J\,:\,a_j = 1}
  \Bigl( \{s : j(s) = j\} \cap \{s : c_j(s) \le H-T\} \Bigr),
\end{equation}
using $f_j = c_j + T \le H$ if and only if $c_j \le H-T$ together with the convention
$(+\infty)+T = +\infty > H$, which is available because $H$ is finite
(Definition~\ref{def:primitives}). Each set in \eqref{eq:D1-cellunion} is Borel by Steps~1 and~3, and
the union is finite by Assumption~\ref{asm:A2p}. Hence $\{D=1\}$ is Borel and $D$ is a measurable
$\{0,1\}$-valued map. Two facts are recorded for later use: at a fixed cutoff policy every ingredient
of \eqref{eq:D1-cellunion} is a function of $s$ alone, so $\{D=1\} \in \sigma(s)$; and Borel
regularity of $s \mapsto B_j(s,d)$ for non-Voice plans was nowhere invoked, the conjunct $a=1$ having
removed those plans from the union.

\emph{Step 5 (exactly one cell, at the level of states).} Put $\mathcal C_F := \{D=1\}$ and
$\mathcal C_P := \{D=0\}$. By Step~4, $D$ is the indicator of a Borel event, so it takes exactly one
of the values $0$ and $1$ at every sample point; therefore $\mathcal C_F \cap \mathcal C_P =
\varnothing$ and $\mathcal C_F \cup \mathcal C_P$ is the whole space. This derives, rather than
asserts, the exclusivity and exhaustiveness recorded in Definition~\ref{def:prices}.

\emph{Step 6 (the assignment is a function of the control-node history, which is part~(a)).} Step~5
partitions the space of states, while part~(a) is a statement about control-node histories. By
Definition~\ref{def:prices} each pooled history $\mathcal H_d^P$ carries the coordinate ``flag landed
by $d$'', so the flag is public by construction, and the same definition supplies the content of the
control-node information set $\mathcal I_H$, which contains the pooled history up to the control
node. By Assumption~\ref{asm:A4} the filing lands exactly at $f = c+T$ and only through the
disclosure node of Section~\ref{sec:timing}, so the coordinate ``flag landed by $H$'' equals
$\ind\{f \le H\}$; combining this with $D = \ind\{a=1,\ c<\infty,\ f \le H\}$ of
Definition~\ref{def:plans} and with Assumption~\ref{asm:A4}'s restriction that only Voice plans
cross, that coordinate equals $D$. Hence $D$ is measurable with respect to the control-node history,
and the map from a history to its cell is single-valued and onto
$\{\mathcal C_F,\mathcal C_P\}$. That is part~(a). The public-flag bridge is what part~(a) turns on:
were the flag coordinate absent from the public history, $D$ would remain a measurable function of
the state by Step~4 while the cell assignment would cease to be a function of what is observed at the
control node.

\emph{Step 7 (part~(b), the forward direction).} Let $j$ be a Voice plan and suppose $f_j \le H$.
Then $c_j = f_j - T \le H-T$ and $c_j < \infty$. By Definition~\ref{def:plans} and
Assumption~\ref{asm:A4}, $c_j$ is the first date at which the path reaches $\tau$, so
$B_j(s,c_j) \ge \tau$. By clause (TR-ii) of Assumption~\ref{asm:TR}, $\partial_d B_j \ge 0$ for
Voice plans, and $c_j \le H-T$, so $B_j(s,H-T) \ge B_j(s,c_j) \ge \tau$.

\emph{Step 8 (part~(b), the converse).} Suppose $B_j(s,H-T) \ge \tau$. Since
$T \in \{1,\dots,H\}$ by Definition~\ref{def:primitives}, the date $H-T$ lies in $\{0,\dots,H-1\}$
and is a calendar date of the model. It therefore belongs to
$\{d \in \{0,\dots,H\} : B_j(s,d) \ge \tau\}$, which is thus nonempty, so its minimum satisfies
$c_j \le H-T < \infty$ by Step~3 and $f_j = c_j+T \le H$. Steps~7 and~8 together give the equivalence
of part~(b) for every Voice plan. The converse direction used no monotonicity in the calendar date;
only the forward direction did. Nor does either direction require a strict crossing: a path sitting
exactly at $\tau$ from its first hitting date through $H-T$ satisfies both sides, the whole argument
running on weak inequalities.

\emph{Step 9 (why the core carries $b_0 < \tau$).} By Definition~\ref{def:plans},
$B_j(s,-1) = b_0$ for every plan, while Hold paths are constant and Exit paths weakly decreasing in
the calendar date by clause (TR-ii) of Assumption~\ref{asm:TR}. If $b_0 \ge \tau$ then a Hold plan
has $B_j(s,0) = b_0 \ge \tau$ and hence $c_j = 0 < \infty$ at a plan with $a_j = 0$, contradicting
Assumption~\ref{asm:A4}'s restriction that only Voice plans cross. Assumption~\ref{asm:A4} and the
path table of clause (TR-ii) are therefore jointly satisfiable only under the maintained restriction
$b_0 < \tau$ of clause (TR-i), which is why a pre-existing crossing is placed outside the core model
rather than treated as a special case.

\emph{Step 10 (a flagged history determines the objects of part~(c)).} On a flagged history the
filing date $f$ is observed, being the first date at which the ``flag landed by $d$'' coordinate of
Definition~\ref{def:prices} equals one (Step~6). By Assumption~\ref{asm:A4} the filing lands exactly
at $c+T$, and $T$ is a known policy parameter, so $c = f-T$ is identified from the history and with
it the baseline date $c^- = c-1$ entering $R_d$ and $R$. By Assumption~\ref{asm:A4} the filing
truthfully reveals the stake, so the reported $B^F$ is $B_j(s,f_j)$ as defined in
Definition~\ref{def:plans}; by the timing of Section~\ref{sec:timing} the flagged order $Q^F$ is
submitted and priced and so belongs to the control-node history. Measurability of these two objects
is not inherited from finiteness and has to be argued, because clause (TR-ii) of
Assumption~\ref{asm:TR} makes stake levels continuum-valued: the finiteness of
Assumption~\ref{asm:A2p} covers the plan menu, the image of $\Gamma$, the noise support and the
calendar, not the stake level. The argument is that $f_j$ is finite-valued and measurable by
Steps~3--8, so that
\begin{equation}
\label{eq:D1-BFmeas}
B_j\bigl(s,f_j(s)\bigr) \;=\; \sum_{d=0}^{H-T} \ind\{f_j(s) = d+T\}\;B_j(s,d)
\end{equation}
is a finite sum of measurable functions, each measurable by Step~3, and likewise
$Q_j^F = b_j^*(s) - B_j^F$.

\emph{Step 11 (every price in part~(c) is a well-defined finite real number).} For
$d \in \{0,\dots,H\}$ the pooled price $P_d^P$ is the fixed point of the pricing map at
$\mathcal H_d^P$, unique by Assumption~\ref{asm:A5}. That map is well defined only if the joint law
of the observed flow and the terminal payoff exists, and this is where clause (TR-ii) of
Assumption~\ref{asm:TR} is consumed in its full form: $X_d = q_{jd} + z_d$ with
$q_{jd} = \Gamma\bigl(B_j(s,d)-B_j(s,d-1)\bigr)$ must be a random variable for \emph{every} plan in
the support, Exit plans included, because pooled pricing integrates over every type; $\Gamma$ is
Borel because it is ordered in the increment. Borel regularity of $s \mapsto B_j(s,d)$ is automatic
for Voice and for Hold and is a genuine addition for Exit, and it is needed here and only here.
Finiteness of the prices is the local boundedness and integrability of Assumption~\ref{asm:A2p}. At
$d = -1$ the price is $P_{-1}^P = \Eop[Y]$ by the first convention of
Definition~\ref{def:prices}, finite for the same reason. The flagged price $P^F$ is the fixed point
at the flagged information set, again by Assumption~\ref{asm:A5}. On the dates: a flagged history has
$c \ge 0$, hence $c^- \ge -1$, and $f^- = c+T-1 \ge 0$ because $T \ge 1$. The convention at $d=-1$ is
not decorative. When $T = H$, Step~8 forces $B_j(s,0) \ge \tau$ and hence $c = 0$ on every flagged
history, so $c^- = -1$ there and the entire run-up path is measured from that base. One reading of
Assumption~\ref{asm:A5} is fixed here as well: since $B^F$ is continuum-valued by Step~10, the
flagged information sets are continuum-indexed, so ``unique fixed point'' is to be read as a
measurably selected family of fixed points rather than a finite list, which is what the continuity
clause of Assumption~\ref{asm:A5} supplies.

\emph{Step 12 (the identity).} By the second convention of Definition~\ref{def:prices},
$P_{\ND}(\mathcal H_{f^-}^P) = P_{f^-}^P$: the not-yet-disclosed price is the last pre-filing pooled
price at the same realised order flow. With the definitions $R = P_{f^-}^P - P_{c^-}^P$ and
$J = P^F - P_{\ND}$ of the same definition,
\begin{equation}
\label{eq:D1-telescope}
R + J \;=\; \bigl(P_{f^-}^P - P_{c^-}^P\bigr) + \bigl(P^F - P_{\ND}\bigr)
       \;=\; \bigl(P_{f^-}^P - P_{c^-}^P\bigr) + \bigl(P^F - P_{f^-}^P\bigr)
       \;=\; P^F - P_{c^-}^P ,
\end{equation}
which is \eqref{eq:runup-jump}. The cancellation of $P_{f^-}^P$ is legitimate because all three
prices are finite real numbers by Step~11. The same-order-flow reading of $P_{\ND}$ is load-bearing
and not cosmetic: under a never-disclosed counterfactual reading the middle terms would not cancel
and the identity would carry the residual $P_{f^-}^P - P_{\ND} \ne 0$.

\emph{Step 13 (the run-up path, and what the identity is).} For $d \in \{c-1,\dots,f-1\}$ the run-up
path $R_d = P_d^P - P_{c^-}^P$ is well defined by Step~11 and computable at the control node by
Step~10, with $R_{c-1} = 0$ and $R_{f-1} = R$. Together with Steps~5, 6, 8, 10 and 12 this
establishes every conjunct of the claim. Two limits on what has been shown are worth stating. The
identity \eqref{eq:D1-telescope} is definitional bookkeeping produced by the same-order-flow
definition of $P_{\ND}$, and this step does not claim any sign for $R$, for $R_d$ or for $J$, nor
$\kappa$-invariance of $J$. And no probability restriction entered Steps~4 through~9: the partition
exists even when $\Prb(D=1)$ is zero or one, so Assumption~\ref{asm:A8} is needed for positive cell
mass and never for the structural partition.
\end{proof}

\begin{proof}[Proof of Lemma~\ref{lem:L1} (L1)]
\emph{Step 1 (the kernel is bounded).} By Definition~\ref{def:premium} the engagement-premium kernel
is $h(\mathcal I) = \pi(\mathcal I)p(\mathcal I)$, and by Definition~\ref{def:prices} the engagement
posterior satisfies $\pi \in [0,1]$ while the entry probability of \eqref{eq:entry} satisfies
$p \in (0,1)$. Their product satisfies $0 \le h \le 1$ pointwise.

\emph{Step 2 (the kernel is a random variable).} The posterior
$\pi(\mathcal I_H) = \Prb(a=1 \mid \mathcal I_H)$ is a version of a conditional probability and is
therefore $\mathcal I_H$-measurable; so is $P(\mathcal I_H) = \Eop[Y \mid \mathcal I_H]$, and
Assumption~\ref{asm:A5} pins \emph{the} version up to a null set, which is what makes $h(\mathcal
I_H)$ a single well-defined random variable rather than an equivalence class from which the
decomposition would have to choose. The entry probability is the composition of these with the
continuous map displayed in \eqref{eq:entry}, hence $\mathcal I_H$-measurable, and a product of
measurable functions is measurable. With Step~1 the kernel is integrable.

\emph{Step 3 (the partition and the weight).} By Lemma~\ref{lem:D1}, $\{D=1\}$ and $\{D=0\}$ are
measurable, disjoint and cover the space, so $\Omega = \Prb(D=1)$ of Definition~\ref{def:premium} is
defined and $\Prb(D=0) = 1-\Omega$.

\emph{Step 4 (the pointwise split).} By Step~3, at every sample point exactly one of
$\ind\{D=1\}$ and $\ind\{D=0\}$ equals one and the other zero, so
$h = h\,\ind\{D=1\} + h\,\ind\{D=0\}$ pointwise.

\emph{Step 5 (integration).} Each summand in Step~4 is bounded in absolute value by $1$ by Step~1 and
is therefore integrable on the single probability space that Assumption~\ref{asm:A1} supplies.
Linearity of the integral gives
$\Eop[h] = \Eop\bigl[h\ind\{D=1\}\bigr] + \Eop\bigl[h\ind\{D=0\}\bigr]$.

\emph{Step 6 (the interior case).} Suppose $0 < \Omega < 1$. Both conditioning events then have
strictly positive probability, so the elementary conditional expectations
$\Eop[h \mid D=1] := \Eop[h\ind\{D=1\}]/\Omega$ and
$\Eop[h \mid D=0] := \Eop[h\ind\{D=0\}]/(1-\Omega)$ are defined and finite, Step~1 bounding both by
$1$. This is the defining property of conditioning on an event of positive probability, not a theorem
being leaned on. Substituting into Step~5,
$\Eop[h] = \Omega\,\Eop[h \mid D=1] + (1-\Omega)\,\Eop[h \mid D=0]$.

\emph{Step 7 (scaling to premia).} The premium wedge $\Delta_m$ is a finite strictly positive
constant, by clause (TR-i) of Assumption~\ref{asm:TR} together with the integrability of
Assumption~\ref{asm:A2p}. Multiplying Step~6 by $\Delta_m$ and reading off the definitions
$\Dact = \Delta_m\Eop[h(\mathcal I_H)]$, $M_F = \Delta_m\Eop[h \mid D=1]$ and
$M_P = \Delta_m\Eop[h \mid D=0]$ of Definition~\ref{def:premium} gives \eqref{eq:L1}.

\emph{Step 8 (the degenerate case $\Omega = 1$).} Then $\Prb(D=0) = 0$, and by Step~1
$\bigl|\Eop[h\ind\{D=0\}]\bigr| \le \Prb(D=0) = 0$, so that term of Step~5 vanishes and
$\Eop[h] = \Eop[h\ind\{D=1\}] = \Omega\,\Eop[h \mid D=1] = \Eop[h \mid D=1]$, the middle equality
being Step~6's definition, available because $\Omega = 1 > 0$. Multiplying by $\Delta_m$ gives
$\Dact = M_F$.

\emph{Step 9 (at $\Omega=1$ the pooled average is not merely unavailable but unidentified).} The
statement's clause that the null-cell average is undefined rather than imputed is a
non-identification claim, and it is proved here rather than asserted. The elementary definition of
Step~6 does not apply to $M_P$, since it would divide by $\Prb(D=0)=0$. The route through the
generated $\sigma$-algebra does not rescue it either. Let $\mathcal D := \sigma(D)$, whose atoms are
$\{D=1\}$ and $\{D=0\}$, and let $x$ be any real number. The function
\begin{equation}
\label{eq:L1-version}
\Delta_m\,\Eop[h \mid D=1]\,\ind\{D=1\} \;+\; x\,\ind\{D=0\}
\end{equation}
is $\mathcal D$-measurable, and it integrates to $\Delta_m\Eop[h\ind_A]$ over every $A \in \mathcal
D$, because two such functions with different values of $x$ differ only on the null set $\{D=0\}$.
Every $x$ therefore yields a version of $\Delta_m\Eop[h \mid \mathcal D]$, so no value of $M_P$ is
determined by the law. The correct statement at $\Omega=1$ is Step~8's one-term equation, not
\eqref{eq:L1} evaluated with a term stipulated to be zero.

\emph{Step 10 (the degenerate case $\Omega = 0$).} Mirror Steps~8 and~9 with the cells exchanged:
$\bigl|\Eop[h\ind\{D=1\}]\bigr| \le \Prb(D=1) = 0$, so $\Eop[h] = \Eop[h \mid D=0]$ and
$\Dact = M_P$, while $M_F$ is unidentified by the argument of Step~9.

Two limits on the reach of this identity are worth recording, because both are easy to overshoot.
First, the identity holds for the disclosure partition that Lemma~\ref{lem:D1} constructs and for no
other: Definition~\ref{def:premium} factors $\Omega = \Prb(a=1)\,\omega_a$, so
$\Omega \ne \Prb(a=1)$ whenever the disclosed share of engagements is below one, and averaging the
kernel over $\{a=1\}$ and $\{a=0\}$ instead gives a different and generally unequal decomposition,
the engaged but undisclosed types sitting in the pooled cell. Second, the factorisation
$\mathcal S = (1-\Omega)\mathcal S_P$ of Theorem~\ref{thm:T1}\,(A) does not follow from
\eqref{eq:L1} alone. Differentiating it in $\kappa$ at fixed policies needs three further
ingredients: $\partial_\kappa M_F = 0$, which is Lemma~\ref{lem:L2}; $\partial_\kappa\Omega = 0$,
which is derived in the proof of Lemma~\ref{lem:L2} below and which Theorem~\ref{thm:T1} carries as
its hypothesis (T-7); and $\kappa$-differentiability of $M_P$, which no standing hypothesis of
Section~\ref{sec:hypotheses} supplies and which Theorem~\ref{thm:T1} therefore carries as (T-8).
\end{proof}

\begin{proof}[Proof of Lemma~\ref{lem:L2} (L2)]
Write $\Xi := (v,s,\xi)$ for the triple whose conditional independence is at issue, $z_{0:H}$ for the
noise vector, $\mathcal H_{f^-}^P$ for the pre-filing pooled history on a flagged path, and $\Sfl$
for the flagged tuple $(B^F,Q^F,a{=}1)$ of Definition~\ref{def:prices}, that is, the filing message
$F$ augmented by the flagged order. The functions $u_1,u_2$ below are bounded and measurable and are
local to this proof.

\emph{Step 1 (where $\kappa$ lives).} At fixed cutoff and execution policies, and under the
no-feedback timing of Assumption~\ref{asm:nofeedback}, every policy object --- the selected plan
$j = j(s)$, the path $B_j(s,\cdot)$, the order marks $q_{jd}(s)$, the terminal target $b_j^*(s)$, the
crossing date $c_j(s)$, the filing date $f_j(s)$, the disclosure indicator $D_j(s)$, the filing stake
$B_j^F(s)$ and the flagged order $Q_j^F(s)$ --- is a function of $(s;\tau,T)$ alone.
Assumption~\ref{asm:nofeedback} removes any dependence on realised order flow or prices, and fixed
policies remove any dependence on $\kappa$ through re-optimised cutoffs. Inspecting the primitives of
Definition~\ref{def:primitives}, noise-trading intensity $\kappa$ appears in exactly one place, the
law of the ternary noise mark; the laws of $v$, $\varepsilon$ and $\xi$ and the remaining constants
carry no $\kappa$. Write $Q_\kappa$ for the law of $z_{0:H}$. At fixed policies, then, $\kappa$
enters the model only through $Q_\kappa$.

\emph{Step 2 (the flagged set and its mass are $\kappa$-free).} By Lemma~\ref{lem:D1} the flagged
cell $\mathcal C_F = \{D=1\}$ is measurable, and by Step~4 of its proof it lies in $\sigma(s)$; by
Step~1 it is one fixed Borel subset of the signal line, the same subset at every $\kappa$. By
Assumption~\ref{asm:A1} and Definition~\ref{def:primitives},
$s = v+\varepsilon \sim N(\mu_v,\sigma_v^2+\sigma_\varepsilon^2)$, a law free of $\kappa$ because the
noise enters neither $v$ nor $\varepsilon$. Hence $\Omega = \Prb(s \in \mathcal C_F)$ is the same
number at every $\kappa$. This yields, as a by-product used nowhere below but consumed by
Theorem~\ref{thm:T1} as its hypothesis (T-7), that $\partial_\kappa\Omega = 0$ at fixed policies.

\emph{Step 3 (the flagged tuple pins the signal, with a measurable inverse).} On $\mathcal C_F$ the
engagement flag is $a=1$, by clause (TR-iii) of Assumption~\ref{asm:TR}, which records
$D = 1 \Rightarrow a = 1$, together with Assumption~\ref{asm:A4}'s restriction that only Voice plans
cross; the third coordinate of $\Sfl$ is therefore uninformative there. By the definition
$Q_j^F = b_j^*(s) - B_j^F$ of Definition~\ref{def:plans},
\begin{equation}
\label{eq:L2-target}
B^F_{j(s)}(s) \;+\; Q^F_{j(s)}(s) \;=\; b^*_{j(s)}(s).
\end{equation}
By Assumption~\ref{asm:A7p} in its on-path injective form A7$'$, the composed terminal target on the
right of \eqref{eq:L2-target} is strictly increasing in the signal on the flagged region and is
therefore injective there. That form is consumed almost surely on the flagged set, which is the
permissive reading and the only coherent one when $B^F$ is continuum-valued, as Step~10 of the proof
of Lemma~\ref{lem:D1} establishes it is: with a continuum-valued flagged tuple no individual tuple
carries positive probability, so on-path injectivity must be read as holding almost surely on a set
of positive probability rather than on a set of atoms. Consequently the first two coordinates of
$\Sfl$ determine $b^*_{j(s)}(s)$, which determines $s$, which at fixed policies determines
$j = j(s)$. Write $\iota_F$ for the resulting inverse. It is strictly increasing on the image of the
composed target and a monotone real function on a Borel set is Borel, so $\iota_F$ is measurable; no
appeal to a measurable-selection theorem is needed here, though injectivity together with
measurability would supply one on standard Borel spaces. Since $\Omega > 0$, the set on which this
holds carries strictly positive probability and the statement is not vacuous. Restricted to
$\mathcal C_F$, therefore, $\sigma(\Sfl) = \sigma(s)$ up to null sets. This is the step at which the
weak identification wording of Assumption~\ref{asm:A7} would not do: it permits two pairs $(j,s)$
with different pooled paths, and \eqref{eq:L2-target} then fails to be injective.

\emph{Step 4 (the pre-filing pooled history is a noise transform indexed by plan and signal).}
Define, for a plan $j$ and a signal $s$,
\begin{equation}
\label{eq:L2-upsilon}
\Upsilon_{j,s}(z_{0:H}) \;:=\;
  \Bigl( \bigl(q_{j0}(s)+z_0,\dots,q_{j,f_j(s)-1}(s)+z_{f_j(s)-1}\bigr);\ \text{flag not landed by }
  f_j(s)-1 \Bigr),
\end{equation}
so that $\mathcal H_{f^-}^P = \Upsilon_{j(s),s}(z_{0:H})$ pointwise on $\mathcal C_F$, by the pooled
history of Definition~\ref{def:prices}. The indexing by $(j,s)$ alone is exactly
Assumption~\ref{asm:nofeedback}: with within-window re-optimisation the marks would depend on
realised flow, hence on past noise, and no such indexing would exist. The map
\eqref{eq:L2-upsilon} is jointly measurable in $(s,z_{0:H})$: the filing date takes finitely many
values because the calendar is finite by Assumption~\ref{asm:A2p}, and on each of the finitely many
level sets of $f_{j(s)}(s)$ --- measurable by Lemma~\ref{lem:D1} --- the map is
$(s,z) \mapsto (q_{j(s)d}(s)+z_d)_{d<t}$, measurable because
$q_{jd} = \Gamma\bigl(B_j(\cdot,d)-B_j(\cdot,d-1)\bigr)$ composes a Borel function, clause (TR-ii) of
Assumption~\ref{asm:TR} giving Voice monotonicity in the signal and only Voice plans reaching
$\mathcal C_F$, with the coarsening $\Gamma$, itself Borel because it is ordered in the increment.
Patching finitely many measurable restrictions gives a measurable map. The second component of
\eqref{eq:L2-upsilon} is a deterministic function of $s$ by Step~1 and carries no information beyond
$s$.

\emph{Step 5 (observed prices enlarge nothing).} By Assumption~\ref{asm:A5} each pooled price is the
unique fixed point of the pricing map at its pooled history and is therefore a function of that
history, and likewise the flagged price is a function of the flagged information set. Adjoining the
value of a function of a $\sigma$-algebra's generator to that $\sigma$-algebra does not enlarge it,
so prices may be ignored when the information content of $\mathcal I_H$ is computed.

\emph{Step 6 (an information sandwich).} On $\mathcal C_F$,
\begin{equation}
\label{eq:L2-sandwich}
\sigma(s) \cap \mathcal C_F \;\subseteq\; \mathcal I_H \;\subseteq\;
  \sigma(s,z_{0:H}) \cap \mathcal C_F .
\end{equation}
For the left inclusion: by the timing of Section~\ref{sec:timing} the filing reveals
$F = (B^F,a{=}1)$ and the flagged order is submitted and priced, so $\Sfl$ belongs to $\mathcal I_H$,
which Definition~\ref{def:prices} sets equal to $\Sfl$ on the flagged cell; by Step~3 the tuple
recovers $s$; and $\mathcal C_F \in \sigma(s)$ by Step~2, so restricting to the flagged cell is
itself an operation within $\sigma(s)$. For the right inclusion: every component of $\mathcal I_H$
--- pooled order flow, flag status, filing content, flagged order and prices --- is a measurable
function of $(s,z_{0:H})$ by Steps~1, 4 and~5. Everything below uses only
\eqref{eq:L2-sandwich}, so the conclusion is robust to the exact content assigned to the
control-node information set, provided it satisfies the sandwich.

\emph{Step 7 (a freezing lemma).} Let $\mathcal G$ be a $\sigma$-algebra with $z_{0:H}$ independent
of $\mathcal G$, let $S$ be $\mathcal G$-measurable, let $w$ be bounded and jointly measurable, and
put $\bar w(x) := \int w(x,\zeta)\,Q_\kappa(\mathrm d\zeta)$. Then
$\Eop[w(S,z_{0:H}) \mid \mathcal G] = \bar w(S)$. For a product form $w(x,\zeta) = u_1(x)u_2(\zeta)$
this is
$\Eop[u_1(S)u_2(z_{0:H}) \mid \mathcal G] = u_1(S)\Eop[u_2(z_{0:H})] = \bar w(S)$, pulling out the
$\mathcal G$-measurable factor and using independence. The product forms generate the product
$\sigma$-algebra and are closed under multiplication; both sides of the asserted equality are linear
in $w$ and stable under bounded monotone limits, so the monotone class theorem extends the equality
to every bounded jointly measurable $w$.

\emph{Step 8 (the conditional independence).} By Step~3, conditioning on $\sigma(\Sfl)$ on the
flagged set is conditioning on $\sigma(s)$, so what has to be shown is
$\mathcal H_{f^-}^P \perp \Xi \mid s$ on $\mathcal C_F$. Given $s$, the middle coordinate of $\Xi$ is
degenerate, so it suffices to show, for bounded measurable $u_1$ on the pre-filing history space and
bounded measurable $u_2$ on $\mathbb R^2$,
\begin{equation}
\label{eq:L2-factor}
\Eop\bigl[u_1(\mathcal H_{f^-}^P)\,u_2(v,\xi) \mid s\bigr]
 \;=\; \Eop\bigl[u_1(\mathcal H_{f^-}^P) \mid s\bigr]\cdot\Eop\bigl[u_2(v,\xi) \mid s\bigr]
 \qquad \text{a.s.\ on } \mathcal C_F,
\end{equation}
which is the standard characterisation of conditional independence given a $\sigma$-algebra. Put
$w(x,\zeta) := u_1\bigl(\Upsilon_{j(x),x}(\zeta)\bigr)$, bounded and jointly measurable by Step~4, so
that $u_1(\mathcal H_{f^-}^P) = w(s,z_{0:H})$. By Assumption~\ref{asm:A1} the noise vector is
independent of $(v,\varepsilon,\xi)$ and hence of $\Xi$, which is a measurable function of that
triple. Apply Step~7 with $\mathcal G = \sigma(v,s,\xi)$ and $S = s$:
$\Eop[w(s,z_{0:H})u_2(v,\xi) \mid v,s,\xi] = u_2(v,\xi)\,\bar w(s)$. Taking
$\Eop[\,\cdot \mid s]$ of both sides and pulling out the $\sigma(s)$-measurable factor $\bar w(s)$
gives $\Eop[u_1u_2 \mid s] = \bar w(s)\,\Eop[u_2 \mid s]$; setting $u_2 \equiv 1$ in the same display
gives $\Eop[u_1 \mid s] = \bar w(s)$, and substituting yields \eqref{eq:L2-factor}. This is the first
conjunct of the statement. The apparent paradox dissolves at exactly this point: unconditionally the
pre-filing pooled history depends strongly on the signal through the order marks, and the
conditioning event removes that dependence by pinning the signal, leaving the noise as the only
remaining randomness in \eqref{eq:L2-upsilon}.

\emph{Step 9 (the flagged posterior).} Let $u_3$ be bounded, $\mathcal I_H$-measurable and supported
on $\mathcal C_F$. By the right inclusion of \eqref{eq:L2-sandwich} and the Doob--Dynkin lemma,
$u_3 = u_5(s,z_{0:H})$ for some bounded jointly measurable $u_5$. Running the computation of Step~8
with $u_5$ in place of $w$ gives $\Eop[u_2u_3] = \Eop\bigl[\bar u_5(s)\,\Eop[u_2 \mid s]\bigr]$,
while the tower property together with $\Eop[u_3 \mid s] = \bar u_5(s)$ from Step~7 gives
$\Eop\bigl[u_3\,\Eop[u_2 \mid s]\bigr] = \Eop\bigl[\bar u_5(s)\,\Eop[u_2 \mid s]\bigr]$. Hence
$\Eop[u_2u_3] = \Eop\bigl[\Eop[u_2 \mid s]\,u_3\bigr]$ for every such $u_3$; since
$\Eop[u_2 \mid s]$ is $\sigma(s)$-measurable and the left inclusion of \eqref{eq:L2-sandwich} places
$\sigma(s) \cap \mathcal C_F$ inside $\mathcal I_H$, it is a version of $\Eop[u_2 \mid \mathcal I_H]$
on the flagged set:
\begin{equation}
\label{eq:L2-posterior}
\Eop[u_2(v,\xi) \mid \mathcal I_H] \;=\; \Eop[u_2(v,\xi) \mid s]
 \qquad \text{a.s.\ on } \mathcal C_F .
\end{equation}
The extension from bounded to integrable $u_2$, needed at Step~11 for $u_2(v,\xi) = v$, follows by
applying \eqref{eq:L2-posterior} to the truncations $\max(-n,\min(n,u_2))$ and passing to the limit
under conditional dominated convergence, the dominating variable being $|u_2|$, integrable because
$v$ is Gaussian by Definition~\ref{def:primitives}. Separately, $\pi(\mathcal I_H) = 1$ on
$\mathcal C_F$, because the flagged cell lies inside $\{a=1\}$ by clause (TR-iii) of
Assumption~\ref{asm:TR} and by Assumption~\ref{asm:A4}'s truthful revelation of purpose, and because
the cell is itself $\mathcal I_H$-measurable by Step~6 of the proof of Lemma~\ref{lem:D1}.

\emph{Step 10 (the posterior is $\kappa$-free).} By \eqref{eq:L2-posterior} the conditional law of
$(v,\xi)$ given $\mathcal I_H$ on the flagged set equals its conditional law given $s$. By
Assumption~\ref{asm:A1} and Definition~\ref{def:primitives} that law is
$v \mid s \sim N\bigl(\mu_v+\beta(s-\mu_v),(1-\beta)\sigma_v^2\bigr)$ with
$\beta = \sigma_v^2/(\sigma_v^2+\sigma_\varepsilon^2)$, and $\xi \mid s \sim N(0,\sigma_\xi^2)$
independent of $v$. None of these depends on $\kappa$, which by Step~1 lives only in $Q_\kappa$ and
enters no coordinate of $\Xi$. With $\pi = 1$ from Step~9, the flagged posterior over $(v,a,\xi)$ is
a $\kappa$-free function of the signal, hence by Step~3 a $\kappa$-free function of $\Sfl$.

\emph{Step 11 (the flagged price).} By clause (TR-iv) of Assumption~\ref{asm:TR} the flagged price
solves the inner fixed point $P = \Eop[Y \mid \mathcal I_H]$ with $Y$ as in \eqref{eq:Y}. On the
flagged cell $a=1$ and $\pi=1$, and $m_0 + \Delta_m = m_1$ by Definition~\ref{def:primitives}. Only
two properties of the bidder-entry rule are used, and they are the two the statement carries as
bookkeeping: (i) $\Prb(\mathsf B = 1 \mid \mathcal I_H) = p(\mathcal I_H)$, and (ii)
$\mathsf B \perp v \mid \mathcal I_H$. The rule \eqref{eq:entry} has both, since under
Assumption~\ref{asm:A1} and Step~10 the bidder's shock satisfies
$\xi \mid \mathcal I_H \sim N(0,\sigma_\xi^2)$ independently of $v$ and the entry indicator is a
function of $\xi$ and of $\mathcal I_H$-measurable quantities; any other rule with these two
properties serves as well. Under them,
\begin{equation}
\label{eq:L2-innerFP}
P \;=\; \bigl(1-p(P)\bigr)\Bigl(\mu_v+\beta(s-\mu_v)+\Delta_V\Bigr) \;+\; p(P)\bigl(P+m_1\bigr),
\qquad
p(P) \;=\; 1-\Phi\!\left(\frac{P+K+m_1-\bar S}{\sigma_\xi}\right),
\end{equation}
where $\Eop[v \mid \mathcal I_H] = \Eop[v \mid s] = \mu_v+\beta(s-\mu_v)$ by
\eqref{eq:L2-posterior}. The data of \eqref{eq:L2-innerFP} are the single number $\Eop[v \mid s]$ and
the parameters of Definition~\ref{def:primitives}, none of which depends on $\kappa$ by Step~10.
Assumption~\ref{asm:A5} makes the fixed point unique, so the solution is a function of the data
alone and equal data give equal solutions; uniqueness is load-bearing rather than decorative here,
since without it an extraneous fixed-point selection could vary with $\kappa$ even though the map
does not. Hence the flagged price is a $\kappa$-free function of the signal and, by Step~3, of
$\Sfl$.

\emph{Step 12 (the entry probability).} By \eqref{eq:entry} with $\pi = 1$ from Step~9, the flagged
entry probability is $p = 1-\Phi\bigl((P^F+K+m_1-\bar S)/\sigma_\xi\bigr)$ on the flagged cell. Every
argument is $\kappa$-free, by Step~11 for the price and by Definition~\ref{def:primitives} for the
constants. Write $p_F(s)$ for it; it is measurable as a composition of the monotone map
$s \mapsto \Eop[v \mid s]$, the fixed-point solution and the normal distribution function, and it
takes values in $(0,1)$.

\emph{Step 13 (the flagged cell average).} By Definition~\ref{def:premium} the kernel is
$h = \pi p$, so $h = p_F(s)$ on the flagged cell by Steps~9 and~12. Since $\Omega > 0$, Step~6 of the
proof of Lemma~\ref{lem:L1} gives
\begin{equation}
\label{eq:L2-MF}
M_F \;=\; \Delta_m\,\Eop[h \mid D=1]
     \;=\; \frac{\Delta_m\,\Eop\bigl[p_F(s)\,\ind\{s \in \mathcal C_F\}\bigr]}{\Omega} .
\end{equation}
The numerator of \eqref{eq:L2-MF} integrates a $\kappa$-free bounded measurable function (Step~12)
over a $\kappa$-free Borel set (Step~2) against the $\kappa$-free law of the signal (Step~2), and the
denominator is $\kappa$-free and strictly positive. Hence $M_F$ is invariant to $\kappa$, which
completes the four invariance conclusions.

Two boundaries of the result are part of it. Nothing above asserts $\kappa$-invariance of the
pre-filing pooled price: that price conditions on order flow whose law is $Q_\kappa$-dependent and
moves with $\kappa$ in general, so with $J = P^F - P_{\ND}$ and $\partial_\kappa P^F = 0$ on the
flagged set, wherever the derivative exists $\partial_\kappa J = -\partial_\kappa P_{f^-}^P$. That
reproduces Definition~\ref{def:prices}'s refusal to claim $\kappa$-invariance of the filing-day jump.
And the result is a fixed-policy statement: if cutoffs or execution paths move with $\kappa$, Step~1
is unavailable and the resulting policy and composition responses are general-equilibrium channels,
which is the territory of Proposition~\ref{prop:C1} rather than of this lemma.
\end{proof}

\begin{proof}[Proof of Lemma~\ref{lem:L3} (L3)]
The order taken here is part~(b) first, since it is a statement about one fixed $\bar\pi$ and
consumes none of the chord restriction, then part~(a), then part~(c), which combines them.
Throughout, $\bar\pi > 0$ and $\kappa$ ranges over an open interval on which the weights of
\eqref{eq:atau} are differentiable.

\emph{Step 1 (what the minimal regularity gives).} By hypothesis $h$ is continuous on $[0,\bar\pi]$
and twice differentiable on the open interval $(0,\bar\pi)$. Twice differentiable on $(0,\bar\pi)$
means that $h'$ exists at every point of that interval and is itself differentiable at every point of
it; a differentiable function is continuous, so $h'$ is continuous on $(0,\bar\pi)$. Nothing is
assumed about $h'$ or $h''$ at either endpoint, and no continuity of $h''$ is assumed anywhere.

\emph{Step 2 (a first difference, and the second difference it produces).} For
$t \in [0,\bar\pi/2]$ set $\delta_h(t) := h\bigl(t+\tfrac{\bar\pi}{2}\bigr) - h(t)$. Evaluating at the
two endpoints of that interval and subtracting,
\begin{equation}
\label{eq:L3-seconddiff}
\delta_h\!\left(\tfrac{\bar\pi}{2}\right) - \delta_h(0)
 = \Bigl[h(\bar\pi)-h\!\left(\tfrac{\bar\pi}{2}\right)\Bigr]
 - \Bigl[h\!\left(\tfrac{\bar\pi}{2}\right)-h(0)\Bigr]
 = h(0) - 2h\!\left(\tfrac{\bar\pi}{2}\right) + h(\bar\pi) = C_h(\bar\pi),
\end{equation}
the last equality being the definition \eqref{eq:chord}. All four evaluations are at points of
$[0,\bar\pi]$, where $h$ is defined by Step~1, so each term is a real number and the cancellation is
arithmetic: the value $h(\bar\pi/2)$ appears once with a minus sign from each bracket, leaving the
coefficient $-2$.

\emph{Step 3 (the regularity of that first difference).} The function $\delta_h$ is continuous on
$[0,\bar\pi/2]$, since for $t$ there both $t$ and $t+\bar\pi/2$ lie in $[0,\bar\pi]$, where $h$ is
continuous. It is differentiable on the open interval $(0,\bar\pi/2)$, since for $t$ there both
$t \in (0,\bar\pi/2) \subset (0,\bar\pi)$ and
$t+\bar\pi/2 \in (\bar\pi/2,\bar\pi) \subset (0,\bar\pi)$ are points at which $h'$ exists by Step~1,
and the derivative of a difference is the difference of derivatives, so
$\delta_h'(t) = h'\bigl(t+\bar\pi/2\bigr) - h'(t)$.

\emph{Step 4 (the first application of the mean value theorem).} Step~3 supplies both hypotheses of
the mean value theorem for $\delta_h$ on $[0,\bar\pi/2]$, so there is $t_1 \in (0,\bar\pi/2)$ with
$\delta_h(\bar\pi/2) - \delta_h(0)
 = \tfrac{\bar\pi}{2}\,\delta_h'(t_1)
 = \tfrac{\bar\pi}{2}\bigl[h'(t_1+\tfrac{\bar\pi}{2}) - h'(t_1)\bigr]$.

\emph{Step 5 (the second application).} Consider the closed interval
$[\,t_1,\ t_1+\bar\pi/2\,]$. Since $0 < t_1 < \bar\pi/2$ by Step~4, its left endpoint satisfies
$t_1 > 0$ and its right endpoint satisfies $t_1+\bar\pi/2 < \bar\pi$, so the interval is contained in
$(0,\bar\pi)$, where $h'$ is continuous and differentiable by Step~1. The mean value theorem applied
to $h'$ on that interval yields $\zeta \in (t_1,\ t_1+\bar\pi/2)$ with
$h'(t_1+\tfrac{\bar\pi}{2}) - h'(t_1) = \tfrac{\bar\pi}{2}h''(\zeta)$.

\emph{Step 6 (part~(b)).} Chaining \eqref{eq:L3-seconddiff} with Steps~4 and~5,
\begin{equation}
\label{eq:L3-mvt}
C_h(\bar\pi) \;=\; \tfrac{\bar\pi}{2}\cdot\tfrac{\bar\pi}{2}\,h''(\zeta)
 \;=\; \tfrac14\,\bar\pi^2\,h''(\zeta),
\qquad \zeta \in \bigl(t_1,\ t_1+\tfrac{\bar\pi}{2}\bigr) \subset (0,\bar\pi),
\end{equation}
which is part~(b). Two features of this derivation are worth recording, because the corollary at
Step~12 does not share them. The identity is exact: no term was discarded and there is no remainder.
And every use of $h$ beyond continuity was at points of the \emph{open} interval $(0,\bar\pi)$, so
neither $h'$ nor $h''$ was invoked at $0$ or at $\bar\pi$; the minimal regularity of Step~1 is
genuinely all that is consumed, and differentiability at the endpoints is not. The statement's gloss
that Darboux's theorem does the rest names the alternative route to \eqref{eq:L3-mvt}, which reaches
the same point from a symmetric two-point average: $h''$ is a derivative on $(0,\bar\pi)$ and
therefore has the intermediate value property, so the arithmetic mean of its values at two points of
a compact subinterval is itself a value of $h''$ at some point between them. Neither route uses
continuity of $h''$, which is why the hypothesis does not carry it.

\emph{Step 7 (differentiating the representation).} Fix $\kappa$. By Assumption~\ref{asm:Atau} the
pooled block's expectation of the kernel is the three-term sum \eqref{eq:atau}. Clause ($\tau$-i) of
that assumption makes the kernel a function of the engagement posterior alone, so the three numbers
$h(0)$, $h(\bar\pi/2)$ and $h(\bar\pi)$ are well defined; clause ($\tau$-ii) fixes the three
evaluation points and $\bar\pi$ itself free of $\kappa$, so those three numbers do not vary with
$\kappa$. A finite sum of products of a constant with a differentiable function of $\kappa$ is
differentiable, with derivative the corresponding sum, so
\begin{equation}
\label{eq:L3-diff}
\partial_\kappa \Eop_\kappa[h] \;=\; A_0'(\kappa)h(0)
 + A_{1/2}'(\kappa)h\!\left(\tfrac{\bar\pi}{2}\right) + A_1'(\kappa)h(\bar\pi).
\end{equation}
Both clauses are load-bearing at exactly this step and neither is free. Without ($\tau$-i) the
right-hand side of \eqref{eq:L3-diff} carries the extra term
$\sum_i A_i(\kappa)\,\partial_\kappa h(\pi_i)$ and part~(a) is false; in the model
$h = \pi\,p$ is a function of two scalars, since \eqref{eq:entry} makes entry depend on the price as
well as on the posterior, so the clause is a restriction and not a reading. Without the $\kappa$-free
$\bar\pi$ half of ($\tau$-ii) it carries the further term
$\bigl[\tfrac12 A_{1/2}h'(\bar\pi/2) + A_1 h'(\bar\pi)\bigr]\partial_\kappa\bar\pi$, which is first
order in $\bar\pi$ against the quadratic vanishing of part~(c), so part~(c) is false without that
half of the clause.

\emph{Step 8 (part~(a)).} Substitute Assumption~\ref{asm:Atau}'s restrictions
$A_0' = A_1' = \Akap$ and $A_{1/2}' = -2\Akap$ into \eqref{eq:L3-diff} and factor out the common
coefficient:
\begin{equation}
\label{eq:L3-parta}
\partial_\kappa \Eop_\kappa[h]
 \;=\; \Akap\Bigl[h(0) - 2h\!\left(\tfrac{\bar\pi}{2}\right) + h(\bar\pi)\Bigr]
 \;=\; \Akap\,C_h(\bar\pi),
\end{equation}
the last equality being \eqref{eq:chord}. This is part~(a), and it is exact. The three restrictions
on the weight derivatives are used at this factorisation and nowhere else: they are precisely the
coefficient pattern $(+1,-2,+1)$ that the second difference carries, which is why the constant of
proportionality is a scalar rather than a triple.

\emph{Step 9 (the same conclusion through the chord gap, and where the symmetry is used).} Let
$\ell_h$ be the affine function on $[0,\bar\pi]$ with $\ell_h(0) = h(0)$ and
$\ell_h(\bar\pi) = h(\bar\pi)$; it is well defined because $\bar\pi > 0$. Since $\ell_h$ and
$h-\ell_h$ sum to $h$ pointwise, $\Eop_\kappa[h] = \Eop_\kappa[\ell_h] + \Eop_\kappa[h-\ell_h]$. The
second term collapses to a single product: $h-\ell_h$ vanishes at $0$ and at $\bar\pi$ by
construction, so under the three-point support of Assumption~\ref{asm:Atau} only the middle atom
survives, and affinity gives $\ell_h(\bar\pi/2) = \tfrac12\bigl(h(0)+h(\bar\pi)\bigr)$, whence
\begin{equation}
\label{eq:L3-gap}
\Eop_\kappa[h-\ell_h]
 \;=\; A_{1/2}(\kappa)\Bigl[h\!\left(\tfrac{\bar\pi}{2}\right)
   - \tfrac12\bigl(h(0)+h(\bar\pi)\bigr)\Bigr]
 \;=\; -\tfrac12\,A_{1/2}(\kappa)\,C_h(\bar\pi).
\end{equation}
Differentiating \eqref{eq:L3-gap} and using $A_{1/2}' = -2\Akap$ returns
$\partial_\kappa\Eop_\kappa[h-\ell_h] = \Akap C_h(\bar\pi)$. The first term contributes nothing:
writing the affine function through its intercept and slope,
$\Eop_\kappa[\ell_h] = \ell_h(0)\,\mathrm{m}(\kappa) + \bigl(h(\bar\pi)-h(0)\bigr)\bar\pi^{-1}
\mathrm{r}(\kappa)$, where $\mathrm{m}$ is the pooled block's total mass and $\mathrm{r}$ its
unnormalised engagement moment, and both are $\kappa$-free by hypothesis, so the derivative vanishes.
The hypothesis $h(0)=0$ makes $\ell_h(0) = 0$ and kills the mass term outright; the moment term needs
its own invariance in either case. This route agrees with Step~8 and displays the mechanism: the
affine part of the kernel contributes no interior motion at all, and the whole motion is carried by
the mass of the single middle atom. It also locates where the three-point symmetry is used, since it
is what collapses the chord-gap sum to one term. For a support with more than one interior atom the
same argument gives $\partial_\kappa\Eop_\kappa[h] = \sum_i A_i'(\kappa)\,(h-\ell_h)(\pi_i)$, a
weighted sum of chord gaps at distinct interior points, which is not a scalar multiple of the second
difference at the midpoint. The two routes are not independent derivations and must not be counted as
two: Step~8 consumes the three weight restrictions, this step consumes one of them together with the
two conservation laws, and Step~15 below shows those input sets are equivalent.

\emph{Step 10 (which normalisation, and the bridge to the pooled sensitivity).} By
Lemma~\ref{lem:D1} the pooled cell is measurable and the two cells are exclusive and exhaustive, so
the pooled block is a genuine cell of the partition and the weights of \eqref{eq:atau} are its atom
masses under either normalisation, conditional on $\{D=0\}$ and summing to one, or unnormalised and
summing to $1-\Omega$. Under the conditional normalisation $\Eop_\kappa[h] = \Eop[h \mid D=0]$, so
Definition~\ref{def:premium} and \eqref{eq:L3-parta} give
\begin{equation}
\label{eq:L3-SP}
\partial_\kappa M_P \;=\; \Delta_m\,\Akap\,C_h(\bar\pi),
\qquad
\mathcal S_P \;=\; \Delta_m\,\lvert \Akap\rvert\,\lvert C_h(\bar\pi)\rvert .
\end{equation}
Under the unnormalised normalisation the weights are $(1-\Omega)$ times the conditional ones, and
$1-\Omega$ is the pooled block's total mass, $\kappa$-free by hypothesis, so the two versions of
\eqref{eq:L3-parta} differ by a $\kappa$-free factor and the identity holds verbatim in both.

\emph{Step 11 (combining parts (a) and (b)).} Applying Step~6 to the kernel and substituting
\eqref{eq:L3-mvt} into \eqref{eq:L3-parta},
\begin{equation}
\label{eq:L3-exact}
\partial_\kappa\Eop_\kappa[h] \;=\; \tfrac14\,\Akap\,\bar\pi^2\,h''(\zeta_{\bar\pi}),
\qquad \zeta_{\bar\pi} \in (0,\bar\pi),
\end{equation}
exactly, at every admissible $\bar\pi$ and $\kappa$. No limit has been taken and no term discarded.

\emph{Step 12 (part~(c), the small-$\bar\pi$ corollary).} Assume in addition the second-order Peano
expansion of the kernel at $0+$, $h(\pi) = h(0)+h'(0)\pi+\tfrac12h''(0)\pi^2+o(\pi^2)$, and that one
and the same kernel serves the whole shrinking family. Substituting the expansion at the three
evaluation points of \eqref{eq:chord}, the constant terms cancel because $1-2+1 = 0$, the linear
terms cancel because $-2\cdot\tfrac12+1 = 0$, the quadratic terms leave
$\bigl(-2\cdot\tfrac18+\tfrac12\bigr)h''(0)\bar\pi^2 = \tfrac14h''(0)\bar\pi^2$, and the three
remainder terms are absorbed into a single $o(\bar\pi^2)$. Hence
$C_h(\bar\pi) = \tfrac14h''(0)\bar\pi^2 + o(\bar\pi^2)$ as $\bar\pi \downarrow 0$. The second clause
is not a technicality: without one kernel along the family, the three evaluations being compared at
different $\bar\pi$ belong to different functions and no limit statement is available at all. A
different sufficient condition, which does not run through Step~6, is that $h''$ extend continuously
to $0$ from the right, since then \eqref{eq:L3-mvt} gives
$\bigl|C_h(\bar\pi)-\tfrac14h''(0)\bar\pi^2\bigr|
= \tfrac14\bar\pi^2\bigl|h''(\zeta_{\bar\pi})-h''(0)\bigr|$ with
$\zeta_{\bar\pi} \to 0$. That route assumes less about the kernel near $0$ and more at $0$; neither
route implies the other, and neither is needed for part~(b).

\emph{Step 13 (the rate at which the interior motion vanishes).} Part~(c) asserts the vanishing of
the interior motion at rate $\bar\pi^2$, and the motion is the product \eqref{eq:L3-parta} rather
than the chord alone. Along the shrinking family,
$\bigl|\partial_\kappa\Eop_\kappa[h]\bigr| = \lvert\Akap\rvert\,\lvert C_h(\bar\pi)\rvert$, so the
rate transfers from the chord to the motion only if $\lvert\Akap\rvert$ is bounded \emph{uniformly in
$\bar\pi$} along the limit, which is the clause the statement carries for the seam at which
Lemma~\ref{lem:L4} consumes this lemma. Boundedness of $\lvert\Akap\rvert$ on $[0,1]$ at a fixed
policy is boundedness in $\kappa$ at one margin and does not deliver this. Under the uniform bound,
\begin{equation}
\label{eq:L3-rate}
\bigl|\partial_\kappa\Eop_\kappa[h]\bigr|
 \;\le\; \tfrac14\,\Bigl(\sup_{\bar\pi}\lvert\Akap\rvert\Bigr)\,\bar\pi^2\,
   \bigl|h''(0)+o(1)\bigr| \;=\; O(\bar\pi^2)
 \qquad (\bar\pi \downarrow 0),
\end{equation}
which is part~(c). The vanishing itself is more robust than the quadratic rate: by
\eqref{eq:L3-parta} the motion is $\Akap C_h(\bar\pi)$, and continuity of $h$ at $0$ alone gives
$C_h(\bar\pi) \to h(0)-2h(0)+h(0) = 0$, so the motion vanishes under continuity at $0$, the uniform
bound and the one-kernel clause, with the two remaining clauses of Step~12 buying the rate rather
than the vanishing.

\emph{Step 14 (the vanishing chord, and why the statement is an ``if'').} Suppose
$C_h(\bar\pi) = 0$. Then \eqref{eq:L3-parta} gives $\partial_\kappa\Eop_\kappa[h] = \Akap\cdot 0 = 0$
at every $\kappa$ in the interval, so the pooled block's expectation is constant there and, by
\eqref{eq:L3-SP}, $\partial_\kappa M_P = 0$ and $\mathcal S_P = 0$. The interior motion is exactly
zero, not small and not merely of indefinite sign. Three consequences belong to the statement rather
than to this case. First, the implication runs one way only. Zero interior motion does not imply a
vanishing chord: take any pooled law satisfying \eqref{eq:atau} whose weights happen to be locally
constant in $\kappa$ on a subinterval, where $\Akap = 0$ and \eqref{eq:L3-parta} returns zero motion
for \emph{every} kernel, including kernels with $C_h(\bar\pi)$ strictly negative. Writing the lemma
with ``if and only if'' would therefore assert something false, which is why it is stated as an
implication and never as an equivalence. Second, a vanishing chord does not make the kernel affine:
by \eqref{eq:L3-mvt} it forces $h''(\zeta_{\bar\pi}) = 0$ at the one point the mean value theorem
produced and at no other, and a kernel strictly concave on part of $[0,\bar\pi]$ and strictly convex
on the rest can have $C_h(\bar\pi) = 0$ exactly. Third, under the hypothesis $h(0)=0$ the vanishing
chord is the testable identity $h(\bar\pi) = 2h(\bar\pi/2)$: the top of the chord is exactly twice
the midpoint.

\emph{Step 15 (what the hypothesis set does, recorded rather than claimed).} The hypotheses of this
lemma are not overdetermined, and it is worth saying why, since two of them look like separate
restrictions and are one. Suppose the pooled block's posterior law is supported on the
$\kappa$-invariant points $\{0,\bar\pi/2,\bar\pi\}$ with weights differentiable in $\kappa$, summing
to a total mass $\mathrm{m}(\kappa)$ and carrying the unnormalised engagement moment
$\mathrm{r}(\kappa) = A_{1/2}(\kappa)\tfrac{\bar\pi}{2}+A_1(\kappa)\bar\pi$. Then
\begin{equation}
\label{eq:L3-equiv}
\bigl[\,\mathrm{m}'(\kappa) = 0 \ \text{ and }\ \mathrm{r}'(\kappa) = 0\,\bigr]
\qquad \Longleftrightarrow \qquad
\bigl[\,A_0' = A_1' \ \text{ and }\ A_{1/2}' = -2A_1'\,\bigr].
\end{equation}
In one direction, $\mathrm{r}' = 0$ reads $A_{1/2}'\tfrac{\bar\pi}{2}+A_1'\bar\pi = 0$, and dividing
by $\bar\pi/2 > 0$ gives $A_{1/2}' = -2A_1'$; substituting that into
$\mathrm{m}' = A_0'+A_{1/2}'+A_1' = 0$ gives $A_0' = 2A_1'-A_1' = A_1'$. In the other, summing the
two restrictions gives $A_1'-2A_1'+A_1' = 0$, so $\mathrm{m}' = 0$, and
$-2A_1'\tfrac{\bar\pi}{2}+A_1'\bar\pi = 0$, so $\mathrm{r}' = 0$. Given the $\kappa$-invariant
three-point support, therefore, the derivative pattern of Assumption~\ref{asm:Atau} and the two
conservation laws assumed here are the same restriction stated twice, and the entire remaining
content of the assumption is its support condition. One further consequence of \eqref{eq:atau} is
recorded because Lemma~\ref{lem:L4} consumes it: reading $\bar\pi$ as the \emph{mean} of the pooled
law rather than as its upper support point is degenerate. With conditional weights summing to one and
mean equal to $\bar\pi$, the moment equation is $\tfrac12 A_{1/2}+A_1 = 1$, which with
$A_0+A_{1/2}+A_1 = 1$ gives $A_0 = A_1-1 \le 0$; nonnegativity then forces $A_0 = 0$, $A_1 = 1$ and
$A_{1/2} = 0$, a point mass at $\bar\pi$ with $\Akap = 0$ and, by \eqref{eq:L3-parta}, zero interior
motion for every kernel. A mean cannot equal the maximum of its own support unless the law is
degenerate, which is why Definition~\ref{def:premium} reads $\bar\pi$ as the upper support point and
excludes the other reading throughout.

The three steps that follow ask what the support condition excludes and what would suffice to
establish it. They are recorded as part of the lemma because they fix the assumption's domain, on
which every conditional statement downstream of it depends; none of them is used in the derivation
of parts~(a)--(c) above.

\emph{Step 16 (a structure that satisfies the support condition, built from the model's own
primitives).} Take a one-round market with a single trading date, engagement binary with
$\Prb(a=1) = \rho = \tfrac12$, the engaging plan's pooled order mark $q_0 = 2\bar z$, and the
non-engaging plan's mark $q_0 = 0$. The engaging mark is admissible because
Definition~\ref{def:plans} sets $q_{jd} = \Gamma\bigl(B_j(s,d)-B_j(s,d-1)\bigr)$ with $\Gamma$ a
finite ordered coarsening and places no ceiling of $\bar z$ on its image, and the non-engaging mark
is the constant Hold path of clause (TR-ii) of Assumption~\ref{asm:TR}. Noise is the ternary mark of
Definition~\ref{def:primitives} and observed flow is $X_0 = q_0+z_0$, so the conditional noise
probabilities allocate as
\begin{equation}
\label{eq:L3-exampleA}
\begin{array}{c|c|ccccc}
a\ (\text{prior}) & q_0 & X_0 = -\bar z & X_0 = 0 & X_0 = \bar z & X_0 = 2\bar z & X_0 = 3\bar z\\
\hline
1\ (\tfrac12) & 2\bar z & \text{---} & \text{---} & \kappa/2 & 1-\kappa & \kappa/2\\
0\ (\tfrac12) & 0 & \kappa/2 & 1-\kappa & \kappa/2 & \text{---} & \text{---}
\end{array}
\end{equation}
Reading \eqref{eq:L3-exampleA} cell by cell: at $X_0 \in \{2\bar z,3\bar z\}$ only the engaging type
contributes, so $\pi = 1$ at joint mass
$\tfrac12(1-\kappa)+\tfrac12\cdot\tfrac{\kappa}{2} = \tfrac{2-\kappa}{4}$; at
$X_0 \in \{-\bar z,0\}$ only the non-engaging type contributes, so $\pi = 0$ at the same joint mass;
and at $X_0 = \bar z$ both contribute, each with $\tfrac12\cdot\tfrac{\kappa}{2}$, so
$\pi = \tfrac12$ at mass $\tfrac{\kappa}{2}$. The support is $\{0,\tfrac12,1\}$, which is
$\{0,\bar\pi/2,\bar\pi\}$ with $\bar\pi = 1$, and every one of the three points is free of $\kappa$:
the two extreme cells are fully revealing at every $\kappa$, and the middle cell's posterior is
$\tfrac12$ because the two types reach $X_0 = \bar z$ through noise realisations of equal probability
$\kappa/2$, which cancels. The weights are $A_1(\kappa) = A_0(\kappa) = \tfrac{2-\kappa}{4}$ and
$A_{1/2}(\kappa) = \tfrac{\kappa}{2}$, summing to one, with $A_0' = A_1' = -\tfrac14$ and
$A_{1/2}' = +\tfrac12 = -2\cdot(-\tfrac14)$: Assumption~\ref{asm:Atau}'s derivative pattern holds
exactly with $\Akap = -\tfrac14$. The moment check of Step~15 confirms it, since
$A_{1/2}\tfrac12+A_1\cdot 1 = \tfrac{\kappa}{4}+\tfrac{2-\kappa}{4} = \tfrac12 = \rho$, free of
$\kappa$. What raising $\kappa$ does here is move mass out of the two revealing end cells into the
single pooling cell, symmetrically and one for one, which is the coefficient pattern $(+1,-2,+1)$
that Step~8 factors. Two features of the construction are load-bearing and neither is a
normalisation. The informed mark must lie strictly outside the reach of the uninformed mark plus
noise, $2\bar z > 0+\bar z$, which is what makes both end cells fully revealing and pins the support
endpoints at $0$ and $1$ at every $\kappa$. And the pre-order engagement share must be exactly
$\tfrac12$, which is what puts the pooling cell's posterior at the midpoint of the chord rather than
elsewhere in it; the word ``symmetric'' in Assumption~\ref{asm:Atau} carries that second requirement.

This example has $\bar\pi = 1$ and so cannot be swept toward zero, which part~(c) requires. A family
in the same class with $\bar\pi$ free is obtained directly: put atoms at $\{0,\bar\pi/2,\bar\pi\}$
with $A_1(\kappa) = A_0(\kappa) = \alpha-c\kappa$ and $A_{1/2}(\kappa) = 1-2\alpha+2c\kappa$, so that
$\Akap = -c$ and the derivative pattern holds by inspection; $\alpha = 0.4$ and $c = 0.3$ keep all
three weights in $[0.1,0.8]$ for $\kappa \in [0,1]$. This is a genuine information structure and not
a list of numbers. For a binary state with prior
$\rho := A_{1/2}\tfrac{\bar\pi}{2}+A_1\bar\pi$, which Step~15 shows is $\kappa$-free and which here
equals $\bar\pi/2$, the likelihoods
\begin{equation}
\label{eq:L3-likelihood}
\Prb(\text{signal } i \mid a=1) = \frac{A_i\pi_i}{\rho},
\qquad
\Prb(\text{signal } i \mid a=0) = \frac{A_i(1-\pi_i)}{1-\rho}
\end{equation}
over the three signals $i$ with $\pi_i \in \{0,\bar\pi/2,\bar\pi\}$ are nonnegative and sum to
$\rho/\rho = 1$ and $(1-\rho)/(1-\rho) = 1$ respectively, and Bayes returns
$\Prb(a=1 \mid i) = A_i\pi_i/\bigl(A_i\pi_i+A_i(1-\pi_i)\bigr) = \pi_i$, which are the posited
posteriors.

\emph{Step 17 (a structure that does not satisfy it, and it is the frozen manuscript's own).} Keep
the single trading date and the ternary noise, but let the engaging plan's mark be $q_0 = +\bar z$
against a non-engaging mark of $0$, with $\Prb(a=1) = \rho \in (0,1)$. Then
$X_0 \in \{-\bar z,0,\bar z,2\bar z\}$ and the cells are
\begin{equation}
\label{eq:L3-exampleB}
\begin{array}{c|l|c}
X_0 & \text{contributing } (q_0,z_0)\text{, unnormalised mass} & \text{posterior}\\
\hline
-\bar z & (0,-\bar z):\ (1-\rho)\tfrac{\kappa}{2}
  & 0\\
0 & (\bar z,-\bar z):\ \rho\tfrac{\kappa}{2};\quad (0,0):\ (1-\rho)(1-\kappa)
  & \pi_-(\kappa) = \dfrac{\rho\kappa/2}{\rho\kappa/2+(1-\rho)(1-\kappa)}\\
\bar z & (\bar z,0):\ \rho(1-\kappa);\quad (0,\bar z):\ (1-\rho)\tfrac{\kappa}{2}
  & \pi_+(\kappa) = \dfrac{\rho(1-\kappa)}{\rho(1-\kappa)+(1-\rho)\kappa/2}\\
2\bar z & (\bar z,\bar z):\ \rho\tfrac{\kappa}{2}
  & 1
\end{array}
\end{equation}
The support condition fails twice over in \eqref{eq:L3-exampleB}. The support has four points rather
than three. And the two interior points move with $\kappa$: the likelihood ratio at $X_0 = 0$ is
$(\kappa/2)/(1-\kappa)$, strictly increasing on $(0,1)$, so $\pi_-$ is strictly increasing, while at
$X_0 = \bar z$ it is $(1-\kappa)/(\kappa/2)$, strictly decreasing, so $\pi_+$ is strictly decreasing.
The consequence for this lemma can be written down exactly. Differentiating
$\Eop_\kappa[h] = \sum_x \Prb(X_0=x)\,h(\pi(x))$ gives
\begin{equation}
\label{eq:L3-twoterm}
\partial_\kappa \Eop_\kappa[h]
 \;=\; \underbrace{\sum_x \bigl[\partial_\kappa\Prb(X_0=x)\bigr]h(\pi(x))}_{\text{the term the
   representation keeps}}
 \;+\; \underbrace{\sum_x \Prb(X_0=x)\,h'(\pi(x))\,\partial_\kappa\pi(x)}_{\text{the term it has no
   room for}},
\end{equation}
and the second sum in \eqref{eq:L3-twoterm} is generically nonzero by the monotonicity just shown.
Even the first sum is not proportional to $C_h(\bar\pi)$: by Step~9 it equals
$\sum_i A_i'(h-\ell_h)(\pi_i)$ over four atoms, a weighted sum of chord gaps at two distinct interior
points, and no single scalar multiple of the second difference at the midpoint reproduces it. This is
not a contrived counterexample. It is the structure the frozen manuscript solves: its no-disclosure
order-flow enumeration has four cells whose two middle posteriors are ratios in which the noise
probabilities $1-\kappa$ and $\kappa/2$ appear in numerator and denominator alike and therefore move
with $\kappa$, and only the two chord ends survive the limits $\kappa \downarrow 0$ and
$\kappa \uparrow 1$. That manuscript's own route to the chord is accordingly not the three-point
representation of \eqref{eq:atau} but the chord-gap identity of Step~9, in which the affine part
contributes a $\kappa$-free constant because the unnormalised engagement moment is pinned and the
interior motion is the motion of the gap between the kernel and its chord. Step~9 transfers to that
structure; Step~8's clean proportionality does not.

\emph{Step 18 (whether the two-round pooled cell is in the satisfying class, and the weakest
sufficient conditions).} This is not settled here and is not claimed either way. Three things can be
stated precisely.

First, Step~16 does not settle it, because the example produces $\bar\pi = 1$ and does so by force
rather than by accident. Within the one-round primitives, non-engaging plans carry marks that are
weakly negative or zero, Hold paths being constant and Exit paths weakly decreasing by clause (TR-ii)
of Assumption~\ref{asm:TR}, while Voice plans carry positive increments; so the largest order-flow
realisation is attainable only by a Voice plan and the top posterior atom is $1$. Any one-round
ternary-noise market with a non-degenerate pooled law therefore has $\bar\pi = 1$, and part~(c)'s
limit $\bar\pi \downarrow 0$ is empty in that class.

Second, the two-round structure is the only place in this model where an endpoint below one could
come from. On the flagged cell the posterior is identically one by Definition~\ref{def:prices}, and
the two-round timing of Section~\ref{sec:timing} removes from the pooled cell exactly the histories
that would carry the revealing top atom into the flagged cell. Whether it does so is a separate
question from whether it could, and the answer at the implemented calibration is that it does not:
Remark~\ref{rem:AtauRecord} records $\bar\pi = 1$ at every non-degenerate node there, because
unflagged Voice types still generate fully revealing pooled order flows. The conjecture that the
two-round timing leaves the pooled cell with a top atom strictly below one is therefore false at that
calibration, and nothing below rests on it.

Third, what would have to be shown can be named, and it is narrower than the assumption looks. By
Step~15 it suffices to establish the support condition alone: that the pooled cell's
engagement-posterior law, at fixed policies, is supported on exactly three points
$0 < \bar\pi/2 < \bar\pi$ with none of them varying with $\kappa$. By Step~16 that decomposes into
two checkable requirements, which are the weakest sufficient conditions this lemma can name:
\begin{enumerate}[label=(S\arabic*),leftmargin=3.2em,itemsep=3pt]
\item every pooled order-flow cell is either fully revealing of $a=0$, or fully revealing of $a=1$ up
  to the pooled cell's own ceiling $\bar\pi$, or a single cell whose posterior is $\bar\pi/2$; and
\item that middle cell's likelihood ratio is free of $\kappa$, which in Step~16 came from the two
  contributing types reaching the cell through noise events of equal probability.
\end{enumerate}
Step~17 fails (S1) and (S2) simultaneously. Whether a two-round plan menu can be built to satisfy
(S1) and (S2) across a whole window of $H$ trading dates --- where the pooled history is the vector
$(X_0,\dots,X_d)$ and the pooled cell is itself carved out by the stake-path event
$\{B_j(s,H-T)<\tau\}$ --- is not determined here, and it is the sense in which the question is open.
At the implemented calibration both conditions are refuted together, and by measurement rather than
by argument: Remark~\ref{rem:AtauRecord} records a support carrying $23$ to $767$ distinct posterior
values with no mass at $\bar\pi/2$ at any node, which is (S1), and interior atoms that move with
$\kappa$ at a two-sided Hausdorff distance reaching $0.4608$ between adjacent-$\kappa$ support sets,
which is (S2). That is applicability evidence at one calibration and it moves no label; the question
is one about the assumption's domain, and a different menu, a different horizon or a different
calibration could still satisfy (S1) and (S2). What the openness costs is stated exactly: it is
load-bearing for this lemma, for leg~3 of Lemma~\ref{lem:L4} and for Part~(B) of
Theorem~\ref{thm:T1} jointly.
\end{proof}

\begin{proof}[Proof of Lemma~\ref{lem:L4} (L4)]
In this proof, $j(\cdot)$ is the plan the frozen cutoff vector selects at each signal,
$\mathcal N := \mathcal C_F(\tau',T)\setminus\mathcal C_F(\tau,T)$ the newly flagged set,
$\nu := \Prb(\mathcal N)$ for its mass, and $\rho_P := 1-\Omega(\tau,T)$ for the pooled mass at the
looser threshold. The three legs are proved in order; nestedness is derived at Step~4 as part of
leg~1 rather than assumed.

\emph{Step 1 (only the clock objects move with the threshold).} At fixed policies the selection map
$j(\cdot)$, the engagement labels and the path family $B_j(s,\cdot)$ are the same objects at $\tau$
and at $\tau'$. Definition~\ref{def:plans} writes the stake path with no threshold argument, whereas
the crossing date, the filing date, the stake at filing and the disclosure indicator all carry one.
In the comparison below, therefore, the only objects that move when the threshold moves are those
four; the path itself is common to the two environments.

\emph{Step 2 (the path is a function of the signal and the date alone).} By the no-feedback timing of
Assumption~\ref{asm:nofeedback} the stake path is a function of the plan, the signal and the date
alone. Composing with Step~1's fixed selection map, $d \mapsto B_{j(s)}(s,d)$ is a function of the
signal and the date: it does not depend on the noise draw, on the realised pooled order flow, on the
pooled prices, or on noise-trading intensity.

\emph{Step 3 (product form of the disclosure indicator, at each threshold).} For every plan, every
signal and each of the two thresholds,
\begin{equation}
\label{eq:L4-product}
D_j(s;\tau,T) \;=\; \ind\{a_j = 1\}\cdot\ind\{B_j(s,H-T) \ge \tau\}.
\end{equation}
If $a_j = 0$ the left side of \eqref{eq:L4-product} is zero by Definition~\ref{def:plans}, whose
first conjunct fails, and the right side is zero through its first factor. If $a_j = 1$ the left side
reduces to $\ind\{c_j<\infty,\ f_j \le H\}$, and the conjunct $c_j < \infty$ is redundant: with
$f_j = c_j+T$ and $T$ finite, $f_j \le H$ implies $c_j \le H-T < \infty$, while the reverse inclusion
is immediate. Lemma~\ref{lem:D1}\,(b), the clock equivalence, then converts $\{f_j \le H\}$ into
$\{B_j(s,H-T) \ge \tau\}$. Two hypotheses put that citation inside its domain.
Assumption~\ref{asm:A4} is what makes the crossing date the first passage and pins the filing to land
exactly at $c_j+T$, which is what Lemma~\ref{lem:D1}\,(b) is an equivalence about; and the
restriction $b_0 < \tau' < \tau$ imposes the core domain at \emph{both} thresholds, without which
Lemma~\ref{lem:D1}\,(b) would be cited off its domain at the tighter one. Combining with Step~1, the
realised indicator $D(s;\tau,T) = \ind\{a_{j(s)}=1\}\cdot\ind\{B_{j(s)}(s,H-T)\ge\tau\}$ is a
function of the signal alone, measurable by Lemma~\ref{lem:D1}\,(a), and free of the noise draw and
of $\kappa$.

\emph{Step 4 (leg~1, the inclusion).} Fix any control-node history with $D_j(s;\tau,T) = 1$. By
\eqref{eq:L4-product} this means $a_j = 1$ and $B_j(s,H-T) \ge \tau$. Since $\tau' < \tau$, the
inequality between real numbers $B_j(s,H-T) \ge \tau > \tau'$ gives $B_j(s,H-T) \ge \tau'$, and
applying \eqref{eq:L4-product} at $\tau'$ to the same plan and signal gives $D_j(s;\tau',T) = 1$. The
\emph{same} number $B_j(s,H-T)$ appears in both comparisons, and this is exactly where fixed policies
and Step~2 are consumed: without path fixity the left-hand side of the comparison at $\tau'$ would be
a different path's stake and the implication would not go through. Since the history was arbitrary,
$\mathcal C_F(\tau,T) \subseteq \mathcal C_F(\tau',T)$, and taking complements in the exclusive and
exhaustive partition of Definition~\ref{def:prices} gives
$\mathcal C_P(\tau',T) \subseteq \mathcal C_P(\tau,T)$: tightening the threshold can only shrink the
pooled cell. Nestedness is a conclusion of this step, not an input to it.

\emph{Step 5 (leg~1, every newly flagged history is generated by a Voice plan).} Every history in
$\mathcal N$ satisfies $D(\tau') = 1$ by the definition of that set, and clause (TR-iii) of
Assumption~\ref{asm:TR} records $D = 1 \Rightarrow a = 1$, which holds because the disclosure
indicator carries $a_j = 1$ as a conjunct. Hence the engagement indicator is identically one on
$\mathcal N$, so every newly flagged history is generated by a Voice plan and, whenever $\nu > 0$,
$\Prb(a=1 \mid \mathcal N) = 1$. This is derived, not assumed; it is consumed at Step~9 and its force
is analysed at Step~11.

\emph{Step 6 (leg~1, the weight comparison).} By Assumption~\ref{asm:A1} the joint law of the
primitives is itself a primitive and does not depend on the threshold, which enters only through the
event $\{D=1\}$. Monotonicity of a probability measure applied to Step~4's inclusion gives
$\Omega(\tau',T) = \Prb\bigl(\mathcal C_F(\tau',T)\bigr) \ge \Prb\bigl(\mathcal C_F(\tau,T)\bigr)
= \Omega(\tau,T)$, completing leg~1.

\emph{Step 7 (the newly flagged set, read on the pooled side).} By Steps~4 and~6 together with
exhaustiveness, $\mathcal N = \mathcal C_P(\tau,T)\setminus\mathcal C_P(\tau',T)$; and since
$\mathcal C_F(\tau,T)$ and $\mathcal N$ are disjoint with union $\mathcal C_F(\tau',T)$, finite
additivity gives $\nu = \Omega(\tau')-\Omega(\tau) \ge 0$, the inequality by Step~6. Read on the
pooled side this is the disjoint union
$\mathcal C_P(\tau,T) = \mathcal C_P(\tau',T)\ \uplus\ \mathcal N$.

\emph{Step 8 (both shares are defined).} The hypothesis $\Omega(\tau',T) < 1$ gives
$\Prb\bigl(\mathcal C_P(\tau',T)\bigr) = 1-\Omega(\tau') > 0$, and by Step~6
$\Omega(\tau) \le \Omega(\tau') < 1$, so $\rho_P = 1-\Omega(\tau) > 0$ as well. Both pooled
engagement shares $\bar\pi_{\mathrm{pr}}(\tau) = \Prb\bigl(a=1 \mid D(\tau)=0\bigr)$ and
$\bar\pi_{\mathrm{pr}}(\tau') = \Prb\bigl(a=1 \mid D(\tau')=0\bigr)$ are therefore defined. Were
$\Omega(\tau') = 1$ instead, the pooled cell at the tighter threshold would be null and
$\bar\pi_{\mathrm{pr}}(\tau')$ undefined rather than imputed, which is the treatment
Lemma~\ref{lem:L1} gives a null cell; the comparison would then have no right-hand side rather than a
wrong one.

\emph{Step 9 (the conditional-probability arithmetic).} Apply finite additivity to the event
$\{a=1\}\cap\mathcal C_P(\tau,T)$ along Step~7's disjoint union. The first resulting term is
$\bar\pi_{\mathrm{pr}}(\tau')(\rho_P-\nu)$, using
$\Prb\bigl(\mathcal C_P(\tau',T)\bigr) = \rho_P-\nu$ from Steps~7 and~8; the second is $\nu\cdot 1$
by Step~5, which is the point at which ``every newly flagged history is Voice'' is consumed; and the
left-hand term is $\bar\pi_{\mathrm{pr}}(\tau)\rho_P$. Dividing by $\rho_P > 0$,
\begin{equation}
\label{eq:L4-convex}
\bar\pi_{\mathrm{pr}}(\tau) \;=\; \Bigl(1-\tfrac{\nu}{\rho_P}\Bigr)\,\bar\pi_{\mathrm{pr}}(\tau')
 \;+\; \tfrac{\nu}{\rho_P}\cdot 1 .
\end{equation}

\emph{Step 10 (leg~2, with its exact identity).} By Step~7, $\mathcal N \subseteq \mathcal
C_P(\tau,T)$, so $0 \le \nu \le \rho_P$ and the weight $\nu/\rho_P$ lies in $[0,1]$, the denominator
being strictly positive by Step~8. Equation \eqref{eq:L4-convex} therefore exhibits the share at the
looser threshold as a convex combination of the share at the tighter one and $1$. Rearranging it
gives the exact identity for the difference between the two shares,
\begin{equation}
\label{eq:L4-share}
\bar\pi_{\mathrm{pr}}(\tau) - \bar\pi_{\mathrm{pr}}(\tau')
 \;=\; \frac{\nu}{\rho_P}\Bigl(1-\bar\pi_{\mathrm{pr}}(\tau')\Bigr) \;\ge\; 0,
\end{equation}
the inequality because $\nu \ge 0$, $\rho_P > 0$ and $\bar\pi_{\mathrm{pr}}(\tau') \le 1$, the last
holding because it is a conditional probability. Hence
$\bar\pi_{\mathrm{pr}}(\tau') \le \bar\pi_{\mathrm{pr}}(\tau)$, which is leg~2. Solving
\eqref{eq:L4-convex} for the tighter share when $\rho_P > \nu$ gives the equivalent closed form
$\bar\pi_{\mathrm{pr}}(\tau') = \bigl(\rho_P\bar\pi_{\mathrm{pr}}(\tau)-\nu\bigr)/(\rho_P-\nu)$,
from which the same difference reads $\nu\bigl(1-\bar\pi_{\mathrm{pr}}(\tau)\bigr)/(\rho_P-\nu)$.

\emph{Step 11 (why engagement share one delivers leg~2 unconditionally).} Were Step~5 weakened, so
that the newly flagged set carried some share $\Prb(a=1\mid\mathcal N) \in [0,1]$ rather than exactly
one, the arithmetic of Step~9 would be unchanged except that the last term of \eqref{eq:L4-convex}
would carry that share, and \eqref{eq:L4-share} would become
$\tfrac{\nu}{\rho_P}\bigl(\Prb(a=1\mid\mathcal N)-\bar\pi_{\mathrm{pr}}(\tau')\bigr)$, whose sign is
the sign of the bracket. The conclusion would then be conditional on the newly flagged histories
being at least as engagement-intensive as the pool they leave, which is an assumption about
\emph{which} histories move. Step~5 supplies the value one, the maximum a conditional probability can
take, so the bracket is nonnegative for every admissible value of $\bar\pi_{\mathrm{pr}}(\tau')$ with
no further restriction. Two corollaries: the inequality is an equality exactly when $\nu = 0$ or
$\bar\pi_{\mathrm{pr}}(\tau') = 1$; and no strict inequality is available without assuming
$\nu > 0$, which is not assumed and is not claimed. Finally, legs~1 and~2 hold at every $\kappa$
simultaneously rather than on average, since by Step~3 the indicator is a function of the signal
alone, by Step~1 so is the engagement label, and by Assumption~\ref{asm:A1} the marginal law of the
signal does not depend on $\kappa$; so $\Omega$, $\nu$ and both shares are invariant to $\kappa$ at
fixed policies, and the statement's common $\kappa$ costs nothing.

\emph{Step 12 (leg~3: the pooled sensitivity in product form at each threshold).}
Lemma~\ref{lem:L3} is consumed here by its statement and not through its proof, with the one
exception noted at Step~13. Its part~(a) gives
$\partial_\kappa\Eop_\kappa[h] = \Akap C_h(\bar\pi)$ at a policy at which
Assumption~\ref{asm:Atau} holds, and reading that through Definition~\ref{def:premium}, which sets
$M_P = \Delta_m\Eop[h \mid D=0]$ and $\mathcal S_P = \lvert\partial_\kappa M_P\rvert$, gives
$\partial_\kappa M_P = \Delta_m\Akap C_h(\bar\pi)$ and hence
$\mathcal S_P = \Delta_m\lvert\Akap\rvert\lvert C_h(\bar\pi)\rvert$ at that policy. Two clauses of
Assumption~\ref{asm:Abr} are what make this available at \emph{both} compared thresholds and in the
form the comparison needs. Clause (br-i) carries the
representation \eqref{eq:atau} for the pooled class under $\tau$ and under $\tau'$, with chord
endpoints $\bar\pi(\tau)$ and $\bar\pi(\tau')$ and coefficients $\Akap(\tau)$ and $\Akap(\tau')$;
clause (br-ii) localises all $\kappa$-dependence of $M_P$ in the weights, so that the total
$\kappa$-derivative that Definition~\ref{def:premium} defines $\mathcal S_P$ to be carries no
composition remainder. The gap that (br-ii) closes is real and not bookkeeping: if the support points
or the kernel moved with $\kappa$, the total derivative would carry the extra term
$\bigl[\tfrac12 A_{1/2}h'(\bar\pi/2)+A_1h'(\bar\pi)\bigr]\partial_\kappa\bar\pi$ together with a term
in $\partial_\kappa h$, and Lemma~\ref{lem:L3} bounds neither. Under the two clauses,
\begin{equation}
\label{eq:L4-SPproduct}
\mathcal S_P(\tau,T) = \Delta_m\,\lvert\Akap(\tau)\rvert\,\lvert C_h(\bar\pi(\tau))\rvert,
\qquad
\mathcal S_P(\tau',T) = \Delta_m\,\lvert\Akap(\tau')\rvert\,\lvert C_h(\bar\pi(\tau'))\rvert .
\end{equation}

\emph{Step 13 (leg~3: transfer from the share to the chord endpoint).} Step~10 is a statement about
the pooled prior engagement share, while \eqref{eq:L4-SPproduct} is written in the chord endpoint.
Clause (br-iv) of Assumption~\ref{asm:Abr} is precisely the assumed link: the endpoint is the upper
support point of the pooled posterior law and is a weakly increasing function of the share, the
\emph{same} function at both thresholds. Applying that function to \eqref{eq:L4-share} gives
$\bar\pi(\tau') \le \bar\pi(\tau)$. Without (br-iv), leg~2 says nothing about the endpoint, since
leg~2 moves a share and the chord is drawn to a support point. The identity branch, in which the
endpoint is read as the mean of the pooled law, is excluded as degenerate by (br-iv) itself and, at
the one point where this proof reads a step of Lemma~\ref{lem:L3}'s proof rather than its statement,
by Step~15 there: under that branch the pooled law is a point mass at the endpoint
with $\Akap = 0$, so \eqref{eq:L4-SPproduct} returns
$\mathcal S_P(\tau) = \mathcal S_P(\tau') = 0$ and leg~3 would hold only because both sides vanish.

\emph{Step 14 (leg~3: the chord comparison).} Assumption~\ref{asm:Atau} maintains
$\lvert C_h\rvert$ weakly increasing in the endpoint, and this is used as the maintained orientation
rather than derived. That property orders two values of \emph{one} functional at a single policy,
while the comparison here reads the chord functional at $\tau$ and at $\tau'$; clause (br-v) of
Assumption~\ref{asm:Abr} is what makes those one functional, by requiring the chord functional and
the kernel it is built from to be the same functions of the posterior at both compared thresholds.
The two clauses do different jobs and both are needed: (br-i) fixes the endpoints and the
coefficients, (br-ii) freezes the support points and the kernel along $\kappa$ only, and (br-iv)
speaks about the endpoint map rather than about the kernel. With them and Step~13,
\begin{equation}
\label{eq:L4-chord}
\lvert C_h(\bar\pi(\tau'))\rvert \;\le\; \lvert C_h(\bar\pi(\tau))\rvert .
\end{equation}
Only the magnitude half of the maintained orientation is read here. The sign half $C_h \le 0$ is
consumed at no step of this leg: \eqref{eq:L4-chord} and \eqref{eq:L4-SPproduct} carry absolute
values throughout, and Step~16 multiplies nonnegative numbers.

\emph{Step 15 (leg~3: the coefficient comparison).} Clause (br-iii) of Assumption~\ref{asm:Abr} gives
$\lvert\Akap(\tau')\rvert \le \lvert\Akap(\tau)\rvert$. This clause cannot be dispensed with. The
coefficient is a property of the pooled information structure, and by Step~7 the pooled class at
$\tau'$ is a strictly different collection of histories from the pooled class at $\tau$ whenever
$\nu > 0$, so nothing in the model ties the two coefficients together; Definition~\ref{def:premium}
bounds the coefficient on $[0,1]$ but says nothing about how it moves with the policy.

\emph{Step 16 (leg~3).} All four factors in \eqref{eq:L4-SPproduct} are nonnegative reals and
$\Delta_m > 0$ by clause (TR-i) of Assumption~\ref{asm:TR}. Multiplying the two weak inequalities
\eqref{eq:L4-chord} and Step~15 and reading \eqref{eq:L4-SPproduct} at each threshold,
\begin{equation}
\label{eq:L4-leg3}
\mathcal S_P(\tau',T)
 = \Delta_m\lvert\Akap(\tau')\rvert\,\lvert C_h(\bar\pi(\tau'))\rvert
 \;\le\; \Delta_m\lvert\Akap(\tau)\rvert\,\lvert C_h(\bar\pi(\tau))\rvert
 = \mathcal S_P(\tau,T),
\end{equation}
which is leg~3. It is the only leg that rests on Assumption~\ref{asm:Abr}.

\emph{Step 17 (the vanishing-chord case).} The maintained orientation is weak, so the boundary case
belongs to the statement rather than being assumed away. If $C_h(\bar\pi(\tau)) = 0$ then
\eqref{eq:L4-chord} gives
$0 \le \lvert C_h(\bar\pi(\tau'))\rvert \le \lvert C_h(\bar\pi(\tau))\rvert = 0$, so the chord
vanishes at the tighter threshold too and \eqref{eq:L4-SPproduct} returns
$\mathcal S_P(\tau') = \mathcal S_P(\tau) = 0$. The conclusion then holds with equality, which is the
equality clause of leg~3.

Three limits on the reach of this lemma are part of it. Every statement above is at fixed policies:
if the blockholder re-optimises when the threshold moves --- a signal that selects an aggressive
Voice plan at $\tau$ switching to Hold at $\tau'$ once the earlier certain flag removes the
pooled-round camouflage --- then the engagement label at that signal changes, Step~4's inclusion
fails, and the flagged weight can fall with a tighter threshold. That is the general-equilibrium
cutoff response, which is the subject of Proposition~\ref{prop:C1}, and this lemma must not be quoted
as an equilibrium statement. No strict inequality is claimed at any leg, since $\nu > 0$ is nowhere
assumed. And nothing here concerns the window margin: the window is held fixed throughout, and a
threshold change bundled with a longer window breaks Step~4 outright, since a history crossing
weakly earlier at $\tau'$ may then file after the horizon.
\end{proof}

